\hypertarget{appendix-f77eadcc4a93}{}
\section*{Complex numbers and planar rotations}
\section{Complex numbers and planar rotations}\label{complex-numbers-and-planar-rotations}

\textbf{Curriculum depth:} Introduced · \textbf{First seen:} Phase6 · \textbf{Latest development:} Phase6

\subsection{Definition and mechanism}\label{definition-and-mechanism}

Complex multiplication combines planar rotation and scaling. A unit complex multiplier represents a rotation; its conjugate gives the inverse. A two-component storage tuple is not automatically a rotation.

Multiplication by \(i\) sends \((x,y)\) to \((-y,x)\). Zero collapses the plane rather than representing any rotation.

\subsection{Prerequisites}\label{prerequisites}

\begin{itemize}
\tightlist
\item
  \hyperref[appendix-52051ad550fe]{Scalars fields and real-number models} --- Field operations define which scalar multipliers the vector laws may use.
\end{itemize}

\subsection{Used by}\label{used-by}

Later curriculum sources can deepen this same node; no future coverage is marked complete.

\subsection{Recall and next depth}\label{recall-and-next-depth}

Derive the rotation equations and separate magnitude drift from angle convention. Full production proofs, target-specific integration and learner mastery require further evidence.

\subsection{Source and navigation}\label{source-and-navigation}

\href{Continuous\%20Notes/Computing\%20Foundations\%20-\%20From\%20Bits\%20to\%20C++\%20and\%20Systems\%20APIs\#Rotation\%20first\%20becomes\%20understandable\%20in\%20a\%20plane}{Teaching narrative} · \href{}{Coverage and source gap (01-Foundations-I, appendix-fd090d982a15 § Phase 6 - Parts 501-600)} · \href{}{Evidence and limits (01-Foundations-I, appendix-a103e7846d21 § Phase 6 - Source analysis and verification)} · \hyperref[appendix-4ad47bdd269a]{Dependency Map}
